\documentclass[20pt,landscape,dvips]{foils} 
% add 'draft' option above to exclude image when compiling
\usepackage[french]{babel}
\usepackage[utf8]{inputenc}  
\usepackage{csquotes} 
%\MakeOuterQuote{"}
\frenchspacing
\DecimalMathComma

\usepackage{latexsym}
\usepackage{amsmath,amssymb,amsfonts}
\usepackage{MnSymbol}
\usepackage{url}
\usepackage{graphicx}
%\usepackage[dvipsnames]{xcolor}
\usepackage{hyperref}
\usepackage{alltt}
\usepackage{pifont,manfnt}
\usepackage[dvipsnames,table]{xcolor}
\usepackage{subfig}
%\usepackage{enumerate}
%\usepackage{colortbl}
\usepackage{multirow,hhline}
\usepackage{cclicenses}
\setlength\parindent{0pt}
\hypersetup{colorlinks=true,citecolor=black,urlcolor=black,linkcolor=black}
\usepackage[style=verbose]{biblatex}
\bibliography{refs}

\newcommand{\highlight}[1]{\textcolor{Plum}{\bfseries #1}}
\newcommand{\remark}[1]{%
%\centerline{
\begin{center}
\framebox[.9\textwidth][t]{
\ding{46} 
\parbox[t]{.8\textwidth}{\small #1}}
\end{center}}

\DeclareMathOperator*{\inlaw}{\sim}
\newcommand{\iid}{\inlaw_{\text{i.i.d.}}}
%\newcommand{\iid}{\mathop{\mathrm{diag}}}
\newcommand{\pobs}{p_{\text{obs}}}

\reversemarginpar
\def\mark{\marginpar{\dbend}}

%\newcommand{\bm}[1]{\mbox{\boldmath{$#1$}}}
\renewcommand{\abstractname}{Summary}
\newcommand\bs{\char '134}   %  a backslash character for the \tt font
%\renewcommand\refname{Additional Readings}

% customize header/footer
\rightheader{}
% Note about the copyleft symbol:
% I use a custom reversed and circled "c" char because \textcopyleft in
% the textcomp package does not support sans serif font.
% Also, "c" is shifted horizontally by 1ex to align with the circle.
%\Restriction{\mbox{\raisebox{1.5ex}{\rotatebox{180}{\textcircled{c\kern.1ex}}} 2009}, \url{www.aliquote.org}}
% now I use the CC licence...
% \Restriction{\cc 2016 \VCRevision}
\Restriction{\cc 2016 Module 11 EESPE}

\title{Méthodes psychométriques en qualité de vie}
\author{Christophe Lalanne\\EA 7334 REMES\\ Unité de Méthodologie des critères d’évaluation\\Université Paris-Diderot, Sorbonne Paris-Cité\\}
\date{\includegraphics[height=18ex]{logo.eps}}

%%% This file has been generated by the vc bundle for TeX.
%%% Do not edit this file!
%%%
%%% Define Git specific macros.
\gdef\GITHash{f4328f7906d309e9224e3e5c2f7f36477f44e69f}%
\gdef\GITAbrHash{f4328f7}%
\gdef\GITParentHashes{136657e4877819e72dc9ebb2f209631a2a8d1057}%
\gdef\GITAbrParentHashes{136657e}%
\gdef\GITAuthorName{Christophe Lalanne}%
\gdef\GITAuthorEmail{ch.lalanne@gmail.com}%
\gdef\GITAuthorDate{2016-07-06 10:29:58 +0200}%
\gdef\GITCommitterName{Christophe Lalanne}%
\gdef\GITCommitterEmail{ch.lalanne@gmail.com}%
\gdef\GITCommitterDate{2016-07-06 10:29:58 +0200}%
%%% Define generic version control macros.
\gdef\VCRevision{\GITAbrHash}%
\gdef\VCAuthor{\GITAuthorName}%
\gdef\VCDateRAW{2016-07-06}%
\gdef\VCDateISO{2016-07-06}%
\gdef\VCDateTEX{2016/07/06}%
\gdef\VCTime{10:29:58 +0200}%
\gdef\VCModifiedText{\textcolor{red}{with local modifications!}}%
%%% Assume clean working copy.
\gdef\VCModified{0}%
\gdef\VCRevisionMod{\VCRevision}%


\begin{document}
\LogoOff
\maketitle
\rightfooter{\quad\textsf{\thepage}}



%---------------------------------------------------------------Slide-
\foilhead{Fidélité de mesure}
\begin{itemize}
\item Consistance interne d'une échelle
\item Stabilité test-retest
\item Accord inter-juges
\item Théorie de la généralisabilité
\end{itemize}



% ---------------------------------------------------------------Slide-
\foilhead{Comment évaluer la fidélité de mesure}

Une formulation alternative consiste à se demander quelles sont les sources
potentielles de variation des scores, et donc comment les mesurer et quel est
leur impact lorsque l'on infère des résultats observés sur un échantillon à la
population ?
Des mesures collectées plusieurs fois à partir d'un même instrument peuvent
survenir de plusieurs manière\autocite{Dunn2000} :
\begin{itemize}
\item évaluation répétée de plusieurs sujets par le même évaluateur ;
\item évaluation alternative d'un même individu par plusieurs évaluateurs ;
\item administration répétée d'un même questionnaire ou de formes parallèles ;
\item utilisation de différentes sous-échelles d'un même questionnaire.
\end{itemize}

% ---------------------------------------------------------------Slide-
\foilhead{Outils statistiques}

La fidélité ou précision de mesure peut être quantifiée à l'aide de différentes
techniques :
\begin{itemize}
\item \highlight{décomposition (linéaire) des composantes de variance en TCT} ;
\item modèles d'équations structurelles ;
\item modèles de réponse à l'item.
\end{itemize}

% ---------------------------------------------------------------Slide-
\foilhead{Fidélité de mesure $\neq$ significativité}

La \highlight{significativité statistique} est utilisée pour évaluer la
probabilité ou la vraisemblance de résultats observés sur un échantillon en
référence à un modèle de population sous l'hypothèse nulle ; la
\highlight{significativité pratique ou clinique} reflète le degré de divergence
des résultats observés avec l'hypothèse nulle (tel que mesuré par une mesure de
taille d'effet) -- sous laquelle on ne distingue pas les patients des sujets
contrôles\autocite{Thompson2003}.

Mais ces deux concepts supposent que les scores sur lesquelles les conclusions
reposent sont des indicateurs corrects et précis de la performance ou de l'état
mesuré chez l'individu.

% ---------------------------------------------------------------Slide-
\foilhead{}

% ---------------------------------------------------------------Slide-
\foilhead{}

Fichier de données et scripts R disponibles à l'adresse suivante :\newline
{\centering \url{https://bitbucket.org/chlalanne/eespe11}\par}
\vfill

\raggedleft \scriptsize -- Typeset with \FoilTeX\ (version 2), Revision \VCRevision

\end{document}
